\section*{Capítulo \ref{ch:b1:spherical-astronomy} --- Astronomía esférica}

\begin{solution}{exo:b1:spherical-astronomy:1}
La \omterm{def:b1:spherical-astronomy:horizontal}{altura} del polo es igual a la \omterm{def:b1:spherical-astronomy:geographic}{latitud}: \ang{52}. Una estrella no se pone nunca
si $\delta > \ang{90} - \ang{52} = \ang{38}$, y no sale nunca si
$\delta < \ang{-38}$.
\end{solution}

\begin{solution}{exo:b1:spherical-astronomy:2}
$\qty{6.75}{h}\times\ang{15} = \ang{101.25}$. $H = \mathrm{LST} - \alpha =
\qty{10}{h} - \qty{6}{h}\,\qty{45}{min} = \qty{3}{h}\,\qty{15}{min}$, es decir,
\ang{48.75}. El \omterm{def:b1:spherical-astronomy:equatorial}{ángulo horario} es positivo: la estrella está al oeste del
meridiano, que cruzó hace tres horas y cuarto sidéreas.
\end{solution}

\begin{solution}{exo:b1:spherical-astronomy:3}
$h = \ang{90} - \varphi + \delta_\odot$: \ang{73.44} en el solsticio de junio,
\ang{50} en los equinoccios y \ang{26.56} en el solsticio de diciembre.
\end{solution}

\begin{solution}{exo:b1:spherical-astronomy:4}
1900: divisible por 100 pero no por 400, no es bisiesto. 2000: divisible por 400,
es bisiesto. 2100: no es bisiesto. El año juliano excede al \omterm{def:b1:spherical-astronomy:tropical-year}{año trópico} en
$365.25 - 365.2422 = 0.0078$ días, \num{0.78} días por siglo; el gregoriano, en
$0.0003$ días, \num{0.03} días por siglo.
\end{solution}

\begin{solution}{exo:b1:spherical-astronomy:5}
$\cos\Delta = \sin^2\ang{45} + \cos^2\ang{45}\cos\ang{90} = 0.5$, luego
$\Delta = \ang{60}$ y la distancia es $6371\times\pi/3 = \qty{6672}{km}$. A lo largo
del paralelo, un \omterm{def:b1:spherical-astronomy:spherical-triangle}{círculo menor} de radio $R\cos\ang{45}$, la distancia es
$6371\times\cos\ang{45}\times\pi/2 = \qty{7076}{km}$. El \omterm{def:b1:spherical-astronomy:spherical-triangle}{círculo máximo} es más
corto en unos \qty{404}{km}.
\end{solution}

\begin{solution}{exo:b1:spherical-astronomy:6}
$\sin h = \sin\ang{48}\sin\ang{20} + \cos\ang{48}\cos\ang{20}\cos\ang{30}
= 0.2542 + 0.5446 = 0.7987$, $h = \ang{53.0}$.
$\cos\mathrm{A} = (\sin\ang{20} - \sin\ang{48}\sin\ang{53.0})/(\cos\ang{48}
\cos\ang{53.0}) = (0.3420 - 0.5935)/0.4027 = -0.6246$, $\arccos = \ang{128.7}$;
como la estrella está al oeste del meridiano ($H > 0$), $\mathrm{A} = \ang{360} -
\ang{128.7} = \ang{231.3}$: hacia el suroeste, a \ang{53} sobre el horizonte.
\end{solution}

\begin{solution}{exo:b1:spherical-astronomy:7}
$\cos H_0 = -\tan\ang{60}\tan\ang{23.44} = -1.7321\times0.4335 = -0.7510$,
$H_0 = \ang{138.67}$; la duración del día es $2\times138.67/15 = \qty{18.49}{h}$,
unas \qty{18}{h}~\qty{29}{min}, como indica la curva de \ang{60} N en el solsticio
de junio.
\end{solution}

\begin{solution}{exo:b1:spherical-astronomy:8}
El paso se adelanta \qty{3.93}{min} cada día, así que $30\times3.93 =
\qty{118}{min}$ en 30 días: la estrella cruza hacia las 20:02. Para cruzar a
mediodía debe adelantarse \qty{10}{h} = \qty{600}{min}, lo que lleva
$600/3.93 \approx 153$ días.
\end{solution}

\begin{solution}{exo:b1:spherical-astronomy:9}
(a) En un equinoccio la \omterm{def:b1:spherical-astronomy:horizontal}{altura} al mediodía es \ang{45}: $I = 1361\times\sin\ang{45}
= \qty{962}{W.m^{-2}}$. (b) En el solsticio de junio $h = \ang{68.44}$ y
$(a/r)^2 = (149.598/152.100)^2 = 0.9674$:
$I = 1361\times0.9674\times\sin\ang{68.44} = \qty{1224}{W.m^{-2}}$. El Sol más alto
compensa de sobra la mayor distancia.
\end{solution}

\begin{solution}{exo:b1:spherical-astronomy:10}
$\sin\delta_\odot = \sin\ang{23.44}\sin\ang{60} = 0.3445$,
$\delta_\odot = \ang{20.15}$. $\tan\alpha_\odot = \cos\ang{23.44}\tan\ang{60}
= 1.5890$, $\alpha_\odot = \ang{57.82}$ (en el mismo cuadrante que $\lambda_\odot$).
La \omterm{def:b1:spherical-astronomy:equatorial}{ascensión recta} va por detrás en \ang{2.18}, es decir, $2.18\times4 \approx
\qty{8.7}{min}$ de tiempo. Un Sol que se moviera uniformemente a lo largo de la
\omterm{def:b1:spherical-astronomy:ecliptic}{eclíptica} seguiría moviéndose de forma irregular en \omterm{def:b1:spherical-astronomy:equatorial}{ascensión recta}, que es la
magnitud que fija el \omterm{def:b1:spherical-astronomy:equatorial}{ángulo horario} y, con él, el reloj de sol: detrás del Sol
medio entre un equinoccio y el siguiente solsticio, delante entre un solsticio y el
siguiente equinoccio. Es el término de oblicuidad de la \omterm{def:b1:spherical-astronomy:solar-time}{ecuación del tiempo}.
\end{solution}

\begin{solution}{exo:b1:spherical-astronomy:11}
En la salida $h = 0$, así que la \cref{prop:b1:spherical-astronomy:transform} da
$\sin\delta = \cos\varphi\cos\mathrm{A}_0$; en el solsticio de junio
$\delta = \varepsilon$, luego $\cos\mathrm{A}_0 = \sin\varepsilon/\cos\varphi$.
Para $\mathrm{A}_0 = \ang{45}$: $\cos\varphi = \sin\ang{23.44}/\cos\ang{45} =
0.5626$, $\varphi = \ang{55.8}$ N. Por encima de $\ang{90} - \varepsilon =
\ang{66.56}$, $\sin\varepsilon/\cos\varphi > 1$: no hay \omterm{def:b1:spherical-astronomy:horizontal}{acimut} de salida, porque el
Sol no se pone ese día (sol de medianoche).
\end{solution}

\begin{solution}{exo:b1:spherical-astronomy:12}
(a) Un huso de ángulo $\theta$ es la fracción $\theta/2\pi$ de la esfera:
$\frac{\theta}{2\pi}\times4\pi R^2 = 2\theta R^2$. (b) Los tres pares de husos
tienen área total $2\times(2A + 2B + 2C)R^2$; cubren la esfera una vez, salvo el
triángulo y su opuesto, cubiertos tres veces cada uno, es decir, dos veces más de
lo necesario: $4(A + B + C)R^2 = 4\pi R^2 + 4S$, donde $S$ es el área del triángulo.
Así, $S = R^2(A + B + C - \pi)$. (c) Tres ángulos rectos: $S = R^2(3\pi/2 - \pi) = \pi R^2/2$, una octava parte de la
esfera (un octante). El triángulo de la figura:
$89.9 + 64.9 + 67.4 - 180 = \ang{42.2}$, es decir, \num{0.736} radianes: $S = 0.736\,R^2$,
\qty{5.9}{\percent} de la esfera.
\end{solution}

\begin{solution}{pb:b1:spherical-astronomy:1}
\textbf{1.} Colatitudes $\ang{90} - (\ang{-24.63}) = \ang{114.63}$ y
$\ang{90} - \ang{51.48} = \ang{38.52}$; diferencia de \omterm{def:b1:spherical-astronomy:geographic}{longitud} \ang{70.40}.

\textbf{2.} $\cos\Delta = \cos\ang{114.63}\cos\ang{38.52} +
\sin\ang{114.63}\sin\ang{38.52}\cos\ang{70.40} = -0.3260 + 0.1899 =
-0.1362$.

\textbf{3.} $\Delta = \ang{97.83}$; distancia $6371\times1.7074 =
\qty{10878}{km}$, unos \qty{10900}{km}.

\textbf{4.}
\[
  \cos\mathrm{Az} = \frac{\cos\ang{38.52} - \cos\ang{114.63}\cos\ang{97.83}}
  {\sin\ang{114.63}\sin\ang{97.83}} = \frac{0.7823 - 0.0568}{0.9006} = 0.8056 ,
\]
$\mathrm{Az} = \ang{36.3}$, hacia el nordeste.

\textbf{5.} $\sin\mathrm{Az} = \sin\ang{70.40}\sin\ang{38.52}/\sin\ang{97.83}
= 0.5922$, que da \ang{36.3} o \ang{143.7}; el coseno positivo de la pregunta 4
elige \ang{36.3}.

\textbf{6.} No. Un rumbo constante corta todos los meridianos con el mismo ángulo y
se enrolla hacia un polo; el rumbo del \omterm{def:b1:spherical-astronomy:spherical-triangle}{círculo máximo} cambia de forma continua (solo
vale \ang{36.3} al principio). El avión recorrería un trayecto más largo y llegaría
a otro punto, salvo que corrigiera el rumbo.

\textbf{7.} La \omterm{def:b1:spherical-astronomy:horizontal}{altura} del \omterm{def:b1:spherical-astronomy:celestial-sphere}{polo celeste} sur es igual al valor absoluto de la \omterm{def:b1:spherical-astronomy:geographic}{latitud}:
\ang{24.63}.

\textbf{8.} $h = \ang{90} - |\ang{-24.63} - (\ang{-16.72})| = \ang{82.09}$;
$\delta > \varphi$, así que Sirio pasa \ang{7.91} al norte del \omterm{def:b1:spherical-astronomy:celestial-sphere}{cenit}.

\textbf{9.} Una estrella no se pone nunca si $\delta < -(\ang{90} - \ang{24.63}) =
\ang{-65.37}$. Canopo, con \ang{-52.70}, no es circumpolar: se pone.

\textbf{10.} $\ang{90} - |\ang{-24.63} + \ang{52.70}| = \ang{61.93}$, al sur del
\omterm{def:b1:spherical-astronomy:celestial-sphere}{cenit}.

\textbf{11.} $\cos H_0 = -\tan(\ang{-24.63})\tan(\ang{-16.72}) =
-(0.4584)(0.3004) = -0.1377$; $H_0 = \ang{97.92}$. Sirio permanece sobre el
horizonte $2\times97.92/15 = \qty{13.06}{h}$ de \omterm{def:b1:spherical-astronomy:sidereal-time}{tiempo sidéreo}, unas \qty{13}{h}.

\textbf{12.} $\cos\mathrm{A}_0 = \sin(\ang{-16.72})/\cos\ang{24.63} =
-0.3165$: Sirio sale con \omterm{def:b1:spherical-astronomy:horizontal}{acimut} \ang{108.5}, al sur del este.

\textbf{13.} Cuando $H = 0$, $\mathrm{LST} = \alpha = \qty{6}{h}~\qty{45}{min}~\qty{9}{s}$.

\textbf{14.} $\qty{86400}{s}\times365.2422/366.2422 = \qty{86164.09}{s} =
\qty{23}{h}~\qty{56}{min}~\qty{4}{s}$; el paso se adelanta
\qty{235.9}{s}, \qty{3}{min}~\qty{56}{s}, cada noche.

\textbf{15.} Hace falta $12 + n\times24/365.2422 = 6.7525 + 24$, así que
$n = 18.7525/0.065710 = 285.4$ días. Doscientos ochenta y cinco días después del 20
de marzo es el 30 de diciembre: Sirio cruza el meridiano a medianoche local media
en la noche del 30 al 31 de diciembre de 2026.

\textbf{16.} Paranal está \ang{70.40} al oeste de Greenwich: la hora solar media
local va por detrás en $70.40/15 = \qty{4.693}{h}$, \qty{4}{h}~\qty{41.6}{min}.

\textbf{17.} \omterm{def:b1:spherical-astronomy:solar-time}{Tiempo solar aparente} = \omterm{def:b1:spherical-astronomy:solar-time}{tiempo solar medio} + \omterm{def:b1:spherical-astronomy:solar-time}{ecuación del tiempo} =
$\qty{0}{h} - \qty{3}{min}$: el reloj de sol marca unas 23:57.

\textbf{18.} $7\times\qty{3.93}{min} \approx \qty{27.5}{min}$ antes.

\textbf{19.} $+\ang{23.44}$ en el solsticio de junio, $-\ang{23.44}$ en el solsticio
de diciembre.

\textbf{20.} Junio: $\ang{90} - |\ang{-24.63} - \ang{23.44}| = \ang{41.93}$.
Diciembre: $\ang{90} - |\ang{-24.63} + \ang{23.44}| = \ang{88.81}$; como
$\delta_\odot = \ang{-23.44} > \varphi$, el Sol pasa \ang{1.19} al norte del \omterm{def:b1:spherical-astronomy:celestial-sphere}{cenit}.

\textbf{21.} $I = 1361\times(1/0.9837)^2\times\sin\ang{88.81} = 1361\times
1.0334\times0.9998 = \qty{1406}{W.m^{-2}}$.

\textbf{22.} En junio, el invierno austral. $\cos H_0 = -\tan(\ang{-24.63})
\tan\ang{23.44} = 0.4584\times0.4335 = 0.1988$: el Sol se pone con el \omterm{def:b1:spherical-astronomy:equatorial}{ángulo
horario} $H_0 = \ang{78.53}$.

\textbf{23.} $\cos\mathrm{A}_0 = \sin\ang{23.44}/\cos\ang{24.63} = 0.4376$:
$\mathrm{A}_0 = \ang{64.0}$, al norte del este.

\textbf{24.} Duración del día $2\times78.53/15 = \qty{10.47}{h}$; la noche más larga
del año en Paranal dura $24 - 10.47 = \qty{13.53}{h}$, es decir,
\textbf{\qty[text-series-to-math=true]{13}{h}~\qty[text-series-to-math=true]{32}{min}}.
\end{solution}
